% =============================================================================
%  Section 4.1 — Cài đặt và phát triển hệ thống
% =============================================================================

\section{Cài đặt và phát triển hệ thống}
\label{sec:system_setup}

Hệ thống được phát triển trên máy cá nhân chạy Windows và đóng gói để triển khai trên Google Cloud Run phục vụ giai đoạn kiểm nghiệm với người dùng (Mục~\ref{sec:user_deployment}). Knowledge Graph Chèo được dựng trên Neo4j, gồm bảy vở chèo cổ, 38 nhân vật, 41 diễn viên, 15 trích đoạn và 27 phiên bản biểu diễn; pipeline GraphRAG với cả ba chiến lược tạo sinh Pre/Mid/Post-Generation được cài đặt đầy đủ và có thể chuyển đổi qua cấu hình. Mô hình ngôn ngữ lớn được sử dụng qua API để đảm bảo tính \textit{LLM-agnostic}; một mô hình duy nhất được chọn làm mô hình chính cho toàn bộ thực nghiệm với tham số sinh cố định (temperature bằng 0, top-p bằng 1) để đảm bảo khả năng tái lập.

Ngoài GraphRAG, hai hệ thống khác được triển khai làm tham chiếu cho các thí nghiệm so sánh ở Mục~\ref{sec:experiments}. Hệ thống thứ nhất là \textbf{RAG truyền thống}. Trong đó, tài liệu mô tả nghệ thuật Chèo được chia thành các đoạn văn bản có độ dài cố định, nhúng véc-tơ bằng một mô hình embedding và lưu trong vector store. Khi nhận câu hỏi, hệ thống tính độ tương đồng cosine giữa véc-tơ câu hỏi và toàn bộ véc-tơ đoạn văn, lấy một số đoạn gần nhất làm ngữ cảnh trước khi gửi tới mô hình ngôn ngữ. Hệ thống thứ hai là \textbf{LLM-only}. Trong hệ thống này, câu hỏi được gửi trực tiếp tới mô hình ngôn ngữ kèm một prompt hệ thống ngắn xác định miền nghệ thuật Chèo, không có cơ chế truy xuất. Hệ thống này có tác dụng tách biệt đóng góp của cơ chế truy xuất so với tri thức nội tại của các LLMs hiện nay.

Cả ba hệ thống sử dụng cùng một mô hình ngôn ngữ và cùng tham số tạo sinh để chênh lệch kết quả chỉ phản ánh khác biệt về nguồn tri thức và cơ chế truy xuất. Prompt của RAG truyền thống và LLM-only được thiết kế đơn giản, chỉ yêu cầu trả lời bằng tiếng Việt dựa trên ngữ cảnh được cung cấp nếu có, không áp dụng các ràng buộc cấu trúc như trong ba chiến lược tạo sinh của GraphRAG. Cách thiết kế này giữ cho các thực nghiệm phản ánh đúng đặc trưng gốc, tránh việc kết quả được cải thiện một cách nhân tạo bởi các kỹ thuật prompting nâng cao.
